\documentclass[11pt,a4paper]{article}
\usepackage[margin=2.6cm]{geometry}
\usepackage{amsmath,amssymb,amsthm}
\usepackage{tikz}
\usepackage{booktabs}
\usepackage[hidelinks]{hyperref}

\newtheorem*{paperthm}{\paperthmname}
\newcommand{\paperthmname}{}
\newtheorem{tool}{Tool}
\theoremstyle{remark}
\newtheorem*{example}{Worked example}
\newtheorem*{checkit}{Check it yourself}
\newcommand{\dd}{\mathrm{d}}
\newcommand{\nb}[1]{\texttt{validation/#1}}

\title{The pizzetti proofs, explained\\[4pt]\large Part 3: Crossing the stellar edge (Section 5.3 and Appendix C)}
\author{Companion to the pizzetti paper (Sections 5.1, 5.3 and Appendix C)}
\date{}

\begin{document}
\maketitle

\begin{abstract}
\noindent When the planet overlaps the edge of the star, the interior series of Part~1 no longer
works. The paper then writes the blocked light as a sum of special functions, turns it into a few
short polynomials (``pieces''), and proves an error bound for each piece (Theorem~3, proof in
Appendix~C). This part explains that proof. Most steps need only first-year calculus. A few need
facts from complex analysis, collected in Section~\ref{sec:tools} as five ``tools'', each with the
reason it is true and a numerical illustration. With the tools in hand the proof is mainly careful
bookkeeping: every quantity is bounded by something computable, and the pieces are added up.
Every step is checked in \nb{04\_hypergeometric\_section5.ipynb} and
\nb{05\_contact\_theorem3.ipynb}.
\end{abstract}

\tableofcontents

%% ===================================================================
\section{The plan in plain words}
\label{sec:plan}

During ingress and egress the planet covers part of the stellar edge. The blocked light
$P_\gamma(z,p)=\int(1-r^2)^\gamma\,\dd A$ (over the covered part of the star) still depends only on
$z$ and $p$, but now it has a ``kink'' at the moment the planet's edge touches the star's edge from
inside (second contact), and the interior series cannot reach past it. The paper proceeds in three
steps:
\begin{enumerate}
  \item \textbf{A formula} (Section~5.1): $P_\gamma$ is written as an infinite sum over $j$ of known
        special functions (Gauss hypergeometric functions) of one variable. Keeping $J$ terms gives
        $P_{\gamma,J}$. \emph{Error 1:} the terms after $J$.
  \item \textbf{Pieces} (Section~5.3): the range of $z$ from first contact down to the cut $z_c$ of
        the interior series is split into five to eight pieces. On each piece $P_{\gamma,J}$ is
        replaced by a polynomial of low degree $N$ (with an explicit $y^{\gamma+3/2}$ and
        logarithm factor at second contact). \emph{Error 2:} the polynomial terms after $N$.
  \item \textbf{Bounds} (Appendix~C): error 1 is bounded by a geometric series (Lemma~2); error 2
        by Cauchy's estimate (Lemma~6), which needs an upper bound $M$ for the size of the function
        on a circle in the complex plane (Lemmas~3--5). The degree $N$ of each piece is the smallest
        for which errors 1 and 2 together stay below the tolerance $\tau$.
\end{enumerate}

\begin{example}
Throughout we use the power-2 law $I=1-0.6(1-\mu^{0.6})$, i.e.\ exponents $\gamma=0$ and $0.3$,
with $p=0.1$ and $\tau=10^{-9}$. The interior series reaches $z_c=0.8203$. The limb region
$0.8203<z<1.1$ is covered by six pieces:
\begin{center}
\begin{tabular}{llcc}
\toprule
piece & range of $z$ & degree $N$ & proven bound\\
\midrule
A, first contact & $1.0462$--$1.1000$ & 6 & $4.4\times10^{-10}$\\
MO, middle (in $k$) & $0.9950$--$1.0462$ & 8 & $5.2\times10^{-10}$\\
MO, middle (in $k$) & $0.9462$--$0.9950$ & 8 & $4.7\times10^{-10}$\\
B, second contact, outside & $0.9000$--$0.9462$ & 7 & $5.9\times10^{-10}$\\
C, second contact, inside & $0.8422$--$0.9000$ & 12 & $8.2\times10^{-10}$\\
MI, middle (in $\kappa$) & $0.8203$--$0.8422$ & 4 & $2.6\times10^{-10}$\\
\bottomrule
\end{tabular}
\end{center}
(Second contact is at $z=1-p=0.9$, first contact at $z=1+p=1.1$.)
\end{example}

%% ===================================================================
\section{Five tools from complex analysis}
\label{sec:tools}

The polynomials are truncated power series. To bound what a truncated power series leaves out, it
pays to let the variable be a complex number $u$. Nothing physical happens there; the complex plane
is just where the cleanest bounds live. All five tools are standard (any textbook on complex
analysis, e.g.\ the first chapters of Ahlfors or of Stein \& Shakarchi).

\begin{tool}[A power series converges in a disk]
A power series $f(u)=\sum_nf_nu^n$ converges inside the largest disk around $0$ that contains no
``bad point'' of $f$ (a point where $f$ is infinite, or where a square root or logarithm has its
branch point).
\end{tool}
Example: $1/(1-u)=1+u+u^2+\dots$ converges for $|u|<1$; the bad point is $u=1$. The same holds for
$\sqrt{1-u}$ or $\ln(1-u)$. This is why each piece is kept small: its expansion point must be far
from the bad points of the function, and then the coefficients fall off fast.

\begin{tool}[Cauchy's estimate]
If $f$ is given by its power series on and inside the circle $|u|=\rho$, and $|f(u)|\le M$ on that
circle, then
\[
  |f_n|\le M\,\rho^{-n}\qquad\text{for every }n.
\]
\end{tool}
\emph{Why.} Put $u=\rho e^{i\theta}$ into the series:
$f(\rho e^{i\theta})=\sum_nf_n\rho^ne^{in\theta}$. This is a Fourier series in $\theta$, so its
coefficients are averages over the circle,
\[
  f_n\rho^n=\frac{1}{2\pi}\int_0^{2\pi}f(\rho e^{i\theta})\,e^{-in\theta}\,\dd\theta ,
\]
and an average is never larger than the largest value: $|f_n|\rho^n\le M$.

\emph{What it buys.} If we know $f$ is at most $M$ on a circle of radius $\rho$, then every
coefficient is bounded, and the terms $f_nu^n$ with $|u|\le w<\rho$ are at most $M(w/\rho)^n$: a
geometric series whose tail we can sum. Example: $f=1/(1-u)$ has $|f|\le10$ on $|u|=0.9$, so
$|f_n|\le10\cdot0.9^{-n}$ (true value: $f_n=1$). Crude, but the factor $(w/\rho)^n$ makes it
harmless after a few dozen terms.

\begin{tool}[Square roots and powers of complex numbers]
A non-zero complex number has two square roots. On a disk where a function $g$ has no zeros, one
can choose $\sqrt g$ (and $g^a$ for any real $a$) so that it varies smoothly, matches the usual
positive root on the real axis, and satisfies $|g^a|=|g|^a$. The \emph{principal} root (the one
with positive real part) is such a choice wherever $\operatorname{Re}g>0$.
\end{tool}
The equality $|g^a|=|g|^a$ is what makes bounds easy: to bound $x^{\gamma+1}$ on a circle, bound
$|x|$.

\begin{tool}[Largest and smallest values lie on the boundary]
For a function given by a power series, the largest value of $|f|$ over a closed disk is reached
on its boundary circle. The real part of such a function takes its smallest and largest values on
the boundary as well.
\end{tool}
\emph{Why (for the real part).} The real part $h$ of such a function has zero Laplacian. By the
mean-value expansion of Part~2 with $\Delta h=0$, its average over any small circle equals its value
at the centre. So $h$ cannot have a strict minimum inside the disk: the average around a minimum
would be larger than the minimum itself.

\begin{tool}[Bounding a hypergeometric series]
The Gauss hypergeometric function is the power series ${}_2F_1(a,b;c;u)=\sum_nt_nu^n$ with $t_0=1$
and term ratio
\[
  \frac{t_{n+1}}{t_n}=\frac{(n+a)(n+b)}{(n+c)(n+1)} .
\]
\begin{enumerate}
  \item[(a)] $|{}_2F_1(u)|\le\sum_n|t_n|\,|u|^n$; if all $t_n>0$, this is ${}_2F_1(|u|)$.
  \item[(b)] The term ratio tends to 1, and from some $n$ on it is at most an explicit number
        $\bar r(n)$ that decreases with $n$ (Lemma~5a of the paper). Once $\bar r(n)\,|u|<1$, the
        rest of the sum is smaller than a geometric series. So $\sum|t_n||u|^n$ can be computed
        with a guaranteed upper bound: add terms one by one, then add the geometric tail.
\end{enumerate}
\end{tool}
For (b), the algebra is one line: $(n+a)(n+b)-(n+c)(n+1)=n(a+b-c-1)+ab-c$, so the ratio is
$1+\frac{n(a+b-c-1)+ab-c}{(n+c)(n+1)}$, and each part of the fraction can only shrink as $n$ grows.

\medskip\noindent\textbf{One thing we take from the handbook.} Near argument 1, a hypergeometric
function splits into a smooth part and a part with a non-integer power (and, for half-integer
$\gamma$, a logarithm):
\[
  {}_2F_1(\dots;1-y)=A(y)+y^{m}\bigl[B_l(y)\ln y+B_c(y)\bigr],
\]
where $A$, $B_l$, $B_c$ are again power series with term-ratio rules. These ``connection formulae''
are in every handbook (DLMF \S15.8); the coefficients are tabulated by pizzetti and checked against
arbitrary-precision evaluations in its test suite. This is the only place where the proof uses a
special-function identity without deriving it.

%% ===================================================================
\section{The formula: blocked light as a sum of Gauss functions (Section 5.1)}

\paragraph{Starting point.} By Green's theorem (the 2D divergence theorem, as in Part~2) the blocked
light can be written as an integral along the part of the planet's rim that lies on the star. With
$s=r^2$ running from $a=(z-p)^2$ to $s_1$ (where $s_1=1$ while the planet crosses the edge and
$s_1=(z+p)^2$ while it is inside),
\[
  P_\gamma=-\frac{1}{2(\gamma+1)}\int_a^{s_1}(1-s)^{\gamma+1}\Bigl[1+\frac{p^2-z^2}{s}\Bigr]
  \frac{\dd s}{\sqrt{(s-a)(b-s)}},\qquad b=(z+p)^2
\]
(paper, Eq.~24). Notebook~4 checks this against the defining area integral to better than $10^{-23}$ of the stellar flux.

\paragraph{Change of variable.} Write $x=1-(z-p)^2$ and $s=1-x(1-t)$. Then $s-a=xt$,
$b-s=4zp(1-kt)$ with
\[
  k=\frac{x}{4zp}=\frac{1-(z-p)^2}{4zp},
\]
and $\dd s=x\,\dd t$. The variable $k$ runs from $0$ at first contact to $1$ at second contact;
inside the disk we use $\kappa=1/k$ instead, which runs from 1 at second contact downwards.

\paragraph{The troublesome $1/s$, expanded.} Since $1-s=x(1-t)$ and $0\le x(1-t)\le x<1$ (for
$z>p$),
\[
  \frac1s=\frac{1}{1-x(1-t)}=\sum_{j\ge0}x^j(1-t)^j ,
\]
a geometric series. Integrating term by term, each term is an Euler integral, i.e.\ a Gauss
function of the single variable $k$ (or $\kappa$):
\[
  H_j(k)=\int_0^1t^{-1/2}(1-t)^{\gamma+1+j}(1-kt)^{-1/2}\,\dd t,\qquad
  G_j(\kappa)=\frac1\pi\int_0^1t^{-1/2}(1-t)^{-1/2}(1-\kappa t)^{\gamma+1+j}\,\dd t .
\]
The result (paper, Eqs.~26--27) is
\[
  P_\gamma=-\frac{x^{\gamma+1}\sqrt k}{2(\gamma+1)}\sum_{j\ge0}c_jH_j(k)\quad\text{(across the edge)},\qquad
  P_\gamma=-\frac{\pi x^{\gamma+1}}{2(\gamma+1)}\sum_{j\ge0}c_jG_j(\kappa)\quad\text{(inside)},
\]
with $c_0=1+p^2-z^2$ and $c_j=(p^2-z^2)x^j$ for $j\ge1$.

\paragraph{Lemma 2 (Error 1): the terms after $j=J$.} Look at the integrals defining $H_j$ and $G_j$.
The integrands are positive, and raising $j$ by one multiplies them by $(1-t)$ or $(1-\kappa t)$,
numbers between 0 and 1. So the integrals \emph{decrease} with $j$: $H_j\le H_{J+1}$ and
$\pi G_j\le\pi G_{J+1}\le\pi$ for $j>J$. Since $|c_j|=|z^2-p^2|\,x^j$, the omitted terms are at
most a geometric series in $x$:
\[
  |P_\gamma-P_{\gamma,J}|\le|\text{prefactor}|\cdot|z^2-p^2|\cdot H_{J+1}\cdot\frac{x^{J+1}}{1-x}
  \quad\text{(and similarly inside, with $\pi$ for $H_{J+1}$)}.
\]
On a piece, each factor is largest at one end of the piece (they all move monotonically with $z$),
which gives a single number for the whole piece (paper, Eq.~C10).

\begin{example}
For the first-contact piece A of our example, $J=7$ terms suffice to make error 1 at most
$2.5\times10^{-11}$ of the stellar flux for the exponent $\gamma=0.3$.
\end{example}

\paragraph{Where this breaks down.} The geometric series converges like $x^J$, and
$x=1-(z-p)^2$ approaches 1 when the planet's centre comes close to the stellar centre ($z\to p$).
For large planets the pieces nearest to $z_c$ then need too many terms. That is why proven pieces
exist only for $p\lesssim0.2$ (at $\tau=10^{-9}$); beyond, pizzetti falls back to numerical
quadrature and says so.

\begin{checkit}
\nb{04\_hypergeometric\_section5.ipynb}: the rim integral and the $j$-series against the defining
area integral (agreement better than $10^{-21}$), and Lemma~2 at single points: error/bound between
0.25 and 0.9999, never above 1.
\end{checkit}

%% ===================================================================
\section{The pieces (Sections 5.2--5.3)}

\paragraph{Why pieces.} By Tool~1 a power series converges up to the nearest bad point. For the
blocked light the bad points are the two contacts (where the flux is not smooth) and, inside the
disk, $\kappa=0$. A single series over the whole limb would converge slowly or not at all. Short
series around well-chosen centres converge fast.

\paragraph{The five kinds of piece.}
\begin{itemize}
  \item \textbf{A (first contact, $k\le\tfrac14$)}: $P_{\gamma,J}=k^{\gamma+3/2}\times$(power series
        in $k$). The factor $k^{\gamma+3/2}$ shows how the blocked light starts from zero.
  \item \textbf{MO (two middle pieces, $\tfrac14\le k\le\tfrac34$)}: Taylor series about $k_0=\tfrac38$
        and $\tfrac58$.
  \item \textbf{B and C (second contact, $y=1-k\le\tfrac14$ outside, $y'=1-\kappa\le\tfrac14$ inside)}:
        $P_{\gamma,J}=R(y)+y^{\gamma+3/2}\bigl[S_L(y)\ln y+S_C(y)\bigr]$, three power series. The
        non-smooth part at second contact is written out explicitly; the three series are smooth.
  \item \textbf{MI (middle pieces inside, in $\kappa$)}: Taylor series of half-width at most $0.07$,
        from the C piece down to $z_c$.
\end{itemize}
Each piece needs the coefficients of its power series. They follow by multiplying and composing
known series (for $z$ as a function of $k$, for $x$, for the Gauss functions); this is exact up to
the order computed (64). On the middle pieces the series of the Gauss functions about $0$ are
truncated after a few hundred terms before being re-centred; Section~\ref{sec:eta} accounts for
that.

%% ===================================================================
\section{Bounding error 2: Lemma 6 (Cauchy's estimate with a correction)}

On a piece with variable $u$, $|u|\le w$, write the smooth part as $f(u)=\sum_nf_nu^n$ and use the
polynomial $\sum_{n\le N}f_nu^n$. What is left out is $\sum_{n>N}f_nu^n$. We split it:
\begin{itemize}
  \item terms $N<n\le64$: their coefficients are computed, so add up $|f_n|w^n$ directly;
  \item terms $n>64$: Cauchy's estimate (Tool~2) on a circle of radius $\rho>w$ gives
        $|f_n|w^n\le M(w/\rho)^n$, a geometric tail $M(w/\rho)^{65}/(1-w/\rho)$.
\end{itemize}
This is the paper's Eq.~(32)/(C15):
\[
  \Bigl|\sum_{n>N}f_nu^n\Bigr|\le\sum_{n=N+1}^{64}|f_n|w^n+M\frac{(w/\rho)^{65}}{1-w/\rho}.
\]
At the contacts, the factors $y^{\gamma+3/2}$ and $y^{\gamma+3/2}|\ln y|$ are replaced by their
largest values on the piece (an elementary calculus exercise, Lemma~5d).

\begin{example}
Piece A, $\gamma=0.3$: the circle radius chosen is $\rho=0.87$ against $w=0.25$, and the bound on
the circle is $M\approx0.2$ (in units of the stellar flux). The far tail is then
$0.2\times(0.25/0.87)^{65}/(1-0.29)\approx10^{-36}$: utterly negligible. The bound of the piece is
dominated by the computed terms $N<n\le64$ and by error 1; in total $4.4\times10^{-10}$, while the
true error of the piece is $2.4\times10^{-10}$.
\end{example}

So the only hard part is a \emph{guaranteed} value of $M$, the size of $f$ on the circle. That is
what Lemmas~3--5 provide.

%% ===================================================================
\section{Getting $M$: bounding the function on a circle (Lemmas 3--5)}

$f$ is a product of simple factors: the prefactor ($x^{\gamma+1}$, $\sqrt k$), the $c_j$
($1+p^2-z^2$ and $(p^2-z^2)x^j$), and the Gauss functions. We bound each on the circle and multiply.

\paragraph{Lemma 3: $z$ and $x$ behave well on the disk.} As a function of $k$ the separation is
$z=p(1-2k)+\sqrt{W}$ with $W=1-4p^2k(1-k)$ (paper, Eq.~28). Two things could go wrong on a complex
disk: the square root could jump between its two values, and $z$ (hence $x$) could be zero, which
would spoil the powers $x^{\gamma+1}$. Neither happens:
\begin{itemize}
  \item If $4p^2|k||1-k|<1$ on the disk, then $W$ lies within distance 1 of the number 1, so
        $\operatorname{Re}W>0$, and the principal root is a smooth choice (Tool~3).
  \item $z$ is never zero, because
        $\bigl(\sqrt W+p(1-2k)\bigr)\bigl(\sqrt W-p(1-2k)\bigr)=W-p^2(1-2k)^2=1-p^2\ne0$.
\end{itemize}
Inside the disk, with $\kappa$, the same holds with $V=(\kappa-2p^2)^2+4p^2(1-p^2)$ in place of $W$.
There the condition $\operatorname{Re}V>0$ is checked with Tool~4: $\operatorname{Re}V$ is the real
part of a power series, so its minimum over the disk lies on the circle, where it is a quadratic in
$\cos\theta$ whose minimum is elementary (paper, Eq.~C11). The code checks these conditions for every
circle it uses and rejects the circle if they fail.

\paragraph{Lemma 4: the size of $z$ and $x$ on the circle, from 1024 samples.} Evaluate $z$ at
1024 equally spaced points of the circle and take the largest $|z|$. Between two samples $z$ can be
a little larger. How much? By Cauchy's estimate on a slightly larger circle (radius $\rho_2$), the
slope $|z'|$ on the circle is at most $C/(\rho_2-\rho)$, where $C$ is a crude upper bound of $|z|$
on the larger disk (from the triangle inequality, e.g.\ $|z|\le p(1+2|k|)+\sqrt{1+4p^2|k||1-k|}$).
Every point of the circle is within arc length $\pi\rho/1024$ of a sample, so
\[
  \max_{\text{circle}}|z|\le\max_{\text{samples}}|z|+\frac{\pi\rho}{1024}\cdot\frac{C}{\rho_2-\rho}.
\]
\emph{Example (piece A):} sampling gives $|z|\le1.3233$ and $|x|\le0.4631$ on the circle; the crude
bounds on the larger disk are $1.363$ and $0.594$; the correction term is tiny.

\paragraph{Lemma 5: the Gauss functions.} Three cases, all from Tool~5:
\begin{itemize}
  \item $H_j(k)$ has positive coefficients, so $|H_j(k)|\le H_j(|k|)$, a real number computed with a
        guaranteed upper bound.
  \item $|G_j(\kappa)|\le1$ when $|1-\kappa|\le1$: in its integral, $1-\kappa t$ lies on the straight
        line between $1$ and $1-\kappa$, both of size at most 1, and the weight averages to 1.
  \item At second contact, each of $A$, $B_l$, $B_c$ is bounded by the sum of the sizes of its
        terms, with a geometric tail. For half-integer $\gamma$ the $B_c$ coefficients contain the
        digamma function $\psi$; the elementary inequality $\ln x-1/x\le\psi(x)\le\ln x$ (for
        $x\ge\tfrac16$, from Binet's formula) bounds them by a decreasing function of $n$
        (Lemma~5b and the paragraph after it).
\end{itemize}
Multiplying the bounds of all factors and summing over $j\le J$ gives $M$ (paper, Eq.~C13 and the
proof of Theorem~3).

%% ===================================================================
\section{The middle pieces: one more correction}
\label{sec:eta}

On the middle pieces the coefficients are obtained by taking the series of each Gauss function
about $k=0$ (or $y=0$), cutting it after a few hundred terms, and re-centring it at the middle of the
piece. Re-centring a polynomial is exact, so what we compute are the exact coefficients of a
slightly different function $\tilde f$, with the Gauss functions replaced by their truncations. The
difference $e=f-\tilde f$ is the truncated tail, which Tool~5 bounds on the circle by a number
$\eta$. Lemma~6 then adds $\eta/(1-w/\rho)$ to the error (paper, Eq.~C16). In practice this term is
far below the tolerance; the point is that it is accounted for rather than assumed small.

%% ===================================================================
\section{Putting it together: Theorem 3}

\renewcommand{\paperthmname}{Theorem 3 (as in the paper, in words)}
\begin{paperthm}
For $p<\tfrac12$, on every piece whose disk conditions (Lemma~3) hold, the difference between the
blocked light $P_\gamma$ and the piece's polynomial approximation is at most error~1 (Lemma~2) plus
error~2 (Lemma~6 with $M$ from Lemmas~3--5). The flux error is the sum of these over the exponents
of the limb-darkening law, weighted by $|c_i|$ and divided by the stellar flux.
\end{paperthm}

\paragraph{The chain of reasoning, in one table.}
\begin{center}
\begin{tabular}{p{4.6cm}p{5.4cm}p{3.6cm}}
\toprule
Step & Why it holds & Checked in\\
\midrule
Blocked light as a sum over $j$ & Green's theorem; geometric series; Euler integral & notebook 4\\
Error 1 (terms after $J$) & integrands shrink with $j$; geometric series & notebook 4 (Lemma 2)\\
Polynomial of the piece & exact series arithmetic to order 64 & notebook 5$^*$\\
Error 2, orders $N+1$ to 64 & coefficients computed & notebook 5$^*$\\
Error 2, orders above 64 & Cauchy's estimate (Tool 2) & notebook 5$^*$\\
$M$: geometry factors & Lemma 3 (Tools 3, 4), Lemma 4 (sampling) & notebook 5$^*$\\
$M$: Gauss functions & Lemma 5 (Tool 5, Euler integral, digamma) & notebook 5$^*$\\
Middle pieces & truncation term $\eta$ & notebook 5$^*$\\
\bottomrule
\end{tabular}
\par\smallskip\footnotesize $^*$ Through the final bound of each piece, not lemma by lemma.
\end{center}

\begin{checkit}
\nb{05\_contact\_theorem3.ipynb} evaluates every certified piece on its own, at points spread over
its range including both ends, and compares with the defining area integral. In the full run
(183 pieces, five laws, $p$ from 0.0092 to 0.2, $\tau=10^{-9}$ and $10^{-12}$) no error exceeded
its bound; the largest ratio was 0.98. The bounds are therefore close to the true errors, which
also shows they are not wildly pessimistic.
\end{checkit}

%% ===================================================================
\section{What the proof does not cover}

\begin{itemize}
  \item \textbf{Rounding.} The bounds are theorems about exact arithmetic. The code evaluates them
        in double precision and multiplies them by $1+10^{-6}$; rounding in the coefficients and in
        the evaluation of the polynomials is at the level of $10^{-16}$ of the stellar flux. A
        bound that is certified down to the last bit would need interval arithmetic.
  \item \textbf{Large planets.} For $p\gtrsim0.2$ (at $\tau=10^{-9}$) some pieces cannot be
        certified (error 1 converges too slowly near $z=p$); those samples are computed by
        quadrature, and pizzetti reports which samples are covered by the proof.
  \item \textbf{The connection formulae} of the Gauss functions are taken from the handbook (DLMF
        \S15.8) and tested numerically, not re-derived.
\end{itemize}

%% ===================================================================
\section{Summary}

Crossing the stellar edge, the blocked light is a sum of Gauss functions of one variable. Cutting
that sum after $J$ terms costs at most a geometric series, because the terms shrink. Each of a few
pieces then turns the sum into a short polynomial. What that polynomial leaves out is bounded by
computing the first 64 coefficients and bounding the rest with Cauchy's estimate: the coefficients
of a function are averages over a circle, so they cannot exceed the function's largest value there.
Finding that largest value is the technical core: the geometry is smooth on the circle (Lemma~3),
its size is pinned down by sampling plus a slope bound (Lemma~4), and the Gauss functions are
bounded by their series or integrals (Lemma~5). None of these steps is deep; there are just many of
them, and each is checked numerically.

\end{document}
